% =====================================================================
%  main.tex — Monografia de Engenharia (draft)
%  Classe: memoir com formatação ABNT customizada
%  Compilação: pdflatex + bibtex (documentado em latexmkrc)
%  Obs.: abntex2 não está disponível no ambiente; usar abntex2 se o
%        modelo institucional exigir (ver MONOGRAFIA_PENDENCIAS.md).
% =====================================================================
\documentclass[12pt,a4paper,oneside,openany]{memoir}

% --- Fontes e codificação ---------------------------------------------
\usepackage[T1]{fontenc}
\usepackage[utf8]{inputenc}
\usepackage{lmodern}
\usepackage{textcomp}

% --- Idioma ------------------------------------------------------------
\usepackage[brazilian]{babel}

% --- Geometria ABNT: 3 cm sup./esq., 2 cm inf./dir. --------------------
\setlrmarginsandblock{3cm}{2cm}{*}
\setulmarginsandblock{3cm}{2cm}{*}
\checkandfixthelayout

% --- Espaçamento 1,5 e recuo de parágrafo ------------------------------
\OnehalfSpacing
\setlength{\parindent}{1.25cm}
\setlength{\parskip}{0pt}

% --- Profundidade do sumário -------------------------------------------
\setcounter{tocdepth}{3}
\setcounter{secnumdepth}{3}
\maxtocdepth{subsection}

% --- Matemática ----------------------------------------------------------
\usepackage{amsmath}
\usepackage{amssymb}
\usepackage{mathtools}

% --- Figuras, tabelas e cores --------------------------------------------
\usepackage{graphicx}
\graphicspath{{figures/}}
\usepackage[table]{xcolor}
\usepackage{booktabs}
\usepackage{longtable}
\usepackage{array}
\usepackage{multirow}
\usepackage{float}
\usepackage{rotating}

% --- Código ---------------------------------------------------------------
\usepackage{listings}
\definecolor{codegreen}{rgb}{0.0,0.42,0.24}
\definecolor{codegray}{rgb}{0.5,0.5,0.5}
\definecolor{codeblue}{rgb}{0.0,0.2,0.6}
\definecolor{codebg}{rgb}{0.96,0.96,0.96}
\lstdefinestyle{monografia}{
  basicstyle=\ttfamily\footnotesize,
  backgroundcolor=\color{codebg},
  commentstyle=\color{codegreen},
  keywordstyle=\color{codeblue}\bfseries,
  stringstyle=\color{codegreen},
  frame=single,
  rulecolor=\color{black!25},
  breaklines=true,
  showstringspaces=false,
  tabsize=2,
  columns=fullflexible,
  captionpos=b,
  literate={á}{{\'a}}1 {à}{{\`a}}1 {ã}{{\~a}}1 {é}{{\'e}}1 {ê}{{\^e}}1
           {í}{{\'i}}1 {ó}{{\'o}}1 {õ}{{\~o}}1 {ú}{{\'u}}1 {ç}{{\c{c}}}1
           {Ã}{{\~A}}1 {É}{{\'E}}1 {Í}{{\'I}}1 {Ó}{{\'O}}1 {Ú}{{\'U}}1 {Ç}{{\c{C}}}1
}
\lstset{style=monografia}

% --- Diagramas TikZ --------------------------------------------------------
\usepackage{tikz}
\usetikzlibrary{arrows.meta,positioning,shapes.geometric,calc,fit,backgrounds,decorations.pathreplacing}

% --- Hiperlinks ---------------------------------------------------------------
\usepackage[hyphens]{url}
\usepackage[hidelinks]{hyperref}
\hypersetup{
  pdftitle={ARGUS - Simulador de defeitos em maquinas rotativas},
  pdfauthor={William da Silva Vianna},
  pdfsubject={Monografia de Engenharia},
  colorlinks=true,
  linkcolor=black,
  urlcolor=black,
  citecolor=black
}

% --- Notas de rodapé e floats -----------------------------------------------
\renewcommand{\footnotesize}{\scriptsize}
\setlength{\abovecaptionskip}{6pt}
\setlength{\belowcaptionskip}{6pt}
\captionnamefont{\bfseries\small}
\captiontitlefont{\small}

% --- Listas de abreviaturas e símbolos (simples) ----------------------------
\usepackage{enumitem}

% =====================================================================
%  PRÉ-TEXTUAIS
% =====================================================================
\begin{document}

\frontmatter
% =====================================================================
%  Capa (modelo institucional a validar — pendência)
% =====================================================================
\begin{titlingpage}
\thispagestyle{empty}
\begin{center}

{\large \textsc{[INSTITUI\c{C}\~AO — Instituto Federal Fluminense]}}\\[0.4cm]
{\large \textsc{[Curso — Engenharia de Controle e Automa\c{c}\~ao]}}\\[4.5cm]

{\large WILLIAM DA SILVA VIANNA}\\[3.5cm]

{\Large\bfseries ARGUS --- Simulador de Defeitos em M\'aquinas Rotativas:}\\[0.3cm]
{\Large\bfseries gera\c{c}\~ao de sinais vibracionais sint\'eticos com}\\[0.3cm]
{\Large\bfseries an\'alise espectral FFT e persist\^encia seletiva para}\\[0.3cm]
{\Large\bfseries treinamento e pesquisa em diagn\'ostico de vibra\c{c}\~ao}\\[5.5cm]

{\large [Cidade]}\\[0.4cm]
{\large 2026}

\end{center}
\end{titlingpage}

% =====================================================================
%  Folha de rosto (natureza do trabalho + banca — a validar)
% =====================================================================
\thispagestyle{empty}
\begin{center}

{\large WILLIAM DA SILVA VIANNA}\\[2.5cm]

{\Large\bfseries ARGUS --- Simulador de Defeitos em M\'aquinas Rotativas:}\\[0.3cm]
{\Large\bfseries gera\c{c}\~ao de sinais vibracionais sint\'eticos, an\'alise}\\[0.3cm]
{\Large\bfseries espectral FFT e persist\^encia seletiva de leituras}\\[3.2cm]

\end{center}

\OnehalfSpacing
\noindent\begin{minipage}{0.48\textwidth}
Monografia de conclus\~ao de curso apresentada ao
[Curso de Engenharia de Controle e Automa\c{c}\~ao] da
[Institui\c{c}\~ao], como requisito parcial para obten\c{c}\~ao do
t\'itulo de [T\'itulo]. \`Area de concentra\c{c}\~ao:
[Engenharia de Controle e Automa\c{c}\~ao].
\end{minipage}

\vspace{2.5cm}

\noindent Orientador: [Nome do orientador --- a definir]

\vspace{1.5cm}
\begin{center}
\begin{minipage}{0.6\textwidth}
\begin{center}
\rule{\textwidth}{0.4pt}\\[0.2cm]
Banca examinadora
\end{center}
\begin{itemize}
  \item Prof. Dr. [Nome] --- [Institui\c{c}\~ao]
  \item Prof. Dr. [Nome] --- [Institui\c{c}\~ao]
  \item Prof. Dr. [Nome] --- [Institui\c{c}\~ao]
\end{itemize}
\end{minipage}
\end{center}

\vfill
\begin{center}
{\large [Cidade]}\\[0.4cm]
{\large 2026}
\end{center}
\clearpage

% =====================================================================
%  Resumo
% =====================================================================
\chapter*{RESUMO}
\addcontentsline{toc}{chapter}{RESUMO}

A manuten\c{c}\~ao preditiva de m\'aquinas rotativas depende da aquisi\c{c}\~ao
cont\'inua de sinais de vibra\c{c}\~ao, cujo volume bruto, quando armazenado de
forma indiscriminada, torna-se um passivo operacional e financeiro. Paralelamente,
o treinamento de analistas e de algoritmos de intelig\^encia artificial em
diagn\'ostico de vibra\c{c}\~ao requer sinais de defeitos variados e
controlados, o que \'e invi\'avel em ativos reais. Este trabalho apresenta o
projeto e a implementa\c{c}\~ao do \textbf{ARGUS}, um simulador de defeitos em
m\'aquinas rotativas que gera sinais vibracionais sint\'eticos com assinatura
espectral parametriz\'avel por tipo de defeito e aplica persist\^encia seletiva
de leituras baseada no crit\'erio \(R^3\) (frequ\^encia, amplitude e fase), com
regra de ouro para varia\c{c}\~ao de rota\c{c}\~ao e Paralelep\'ipedo de
Descarte para filtragem de leituras redundantes.

O sistema \'e composto por um back-end em Python e FastAPI, com m\'odulos de
gera\c{c}\~ao de sinal, processamento espectral por FFT, motor de descarte e
persist\^encia em PostgreSQL 16, e por um front-end em React, que apresenta
painel de sinal, espectro FFT com ordens normalizadas \(N\) e rota\c{c}\~ao em
Hz, indicadores de processamento e taxa de descarte e registro de \emph{snapshots}
de defeito. O cat\'alogo contempla treze tipos de defeito com assinaturas
espectrais derivadas da literatura de diagn\'ostico de vibra\c{c}\~ao, incluindo
desbalanceamento, desalinhamento angular e paralelo, folgas, ro\c{c}amento,
instabilidades de filme de \'oleo e falhas de rolamento (BPFO, BPFI, BSF e FTF).

A valida\c{c}\~ao foi conduzida por testes automatizados (63 testes no
back-end e 12 no front-end) e por ensaio de ponta a ponta no navegador, com
verifica\c{c}\~ao das assinaturas espectrais geradas pela API. Os resultados
indicam que o simulador reproduz as assinaturas esperadas e que o motor de
descarte executa a persist\^encia seletiva conforme a hierarquia de decis\~ao
definida, habilitando a cria\c{c}\~ao de bases de dados de treinamento e a
pesquisa em an\'alise de causa raiz sem exposi\c{c}\~ao de ativos reais.

\vspace{0.6cm}
\noindent\textbf{Palavras-chave:} manuten\c{c}\~ao preditiva; an\'alise de
vibra\c{c}\~ao; FFT; gera\c{c}\~ao de sinais sint\'eticos; persist\^encia
seletiva; monitoramento de condi\c{c}\~ao.

% =====================================================================
%  Abstract
% =====================================================================
\chapter*{ABSTRACT}
\addcontentsline{toc}{chapter}{ABSTRACT}

Predictive maintenance of rotating machinery depends on the continuous
acquisition of vibration signals whose raw volume, when stored
indiscriminately, becomes an operational and financial liability. At the same
time, the training of analysts and of artificial-intelligence algorithms in
vibration diagnostics requires varied and controlled defect signals, which is
impractical on real assets. This work presents the design and implementation of
\textbf{ARGUS}, a fault simulator for rotating machinery that generates synthetic
vibration signals with a spectral signature parameterized by fault type and
applies selective persistence of readings based on the \(R^3\) criterion
(frequency, amplitude and phase), with a golden rule for rotational-speed
variation and a Discard Parallelepiped to filter redundant readings.

The system comprises a Python/FastAPI backend with signal-generation,
FFT-spectral-processing, discard-engine and PostgreSQL~16 persistence modules,
and a React frontend presenting a time-signal panel, an FFT spectrum with
normalized orders \(N\) and rotation in Hz, processing and discard-rate
indicators, and fault-snapshot registration. The catalog covers thirteen fault
types with spectral signatures derived from the vibration-diagnostics
literature, including unbalance, angular and parallel misalignment, looseness,
rotor rub, oil-film instabilities and rolling-element bearing faults (BPFO,
BPFI, BSF and FTF).

Validation was carried out through automated tests (63 backend and 12 frontend
tests) and an end-to-end browser trial, with verification of the spectral
signatures produced by the API. The results indicate that the simulator
reproduces the expected signatures and that the discard engine performs
selective persistence according to the defined decision hierarchy, enabling the
creation of training datasets and root-cause-analysis research without exposing
real assets.

\vspace{0.6cm}
\noindent\textbf{Keywords:} predictive maintenance; vibration analysis; FFT;
synthetic signal generation; selective persistence; condition monitoring.


% Listas
\renewcommand{\listfigurename}{LISTA DE FIGURAS}
\renewcommand{\listtablename}{LISTA DE TABELAS}
\renewcommand{\contentsname}{SUM\'ARIO}
\pdfbookmark[0]{\contentsname}{tof}
\tableofcontents*
\clearpage
\pdfbookmark[0]{\listfigurename}{lof}
\listoffigures*
\clearpage
\pdfbookmark[0]{\listtablename}{lot}
\listoftables*
\clearpage

% =====================================================================
%  TEXTUAIS
% =====================================================================
\mainmatter

% =====================================================================
%  Capítulo 1 — Introdução
% =====================================================================
\chapter{Introdução}
\label{cap:introducao}

\section{Contextualização}

A disponibilidade operacional de máquinas rotativas --- motores, bombas,
compressores, turbinas e redutores --- é um fator crítico em plantas
industriais de alta performance. A falha não planejada de um ativo não apenas
interrompe a produção, mas pode gerar danos secundários, riscos de segurança e
custos de reparo muito superiores aos de uma intervenção planejada. Para
mitigar esses riscos, a engenharia de manutenção evoluiu da manutenção
corretiva e preventiva baseada em calendário para a \emph{manutenção
preditiva}, na qual a condição do ativo é monitorada continuamente e as
intervenções são decididas com base em evidências do estado de degradação
\cite{scheffer2004,randall2011}.

Dentre as técnicas de monitoramento de condição, a análise de vibração é uma
das mais difundidas e informativas para máquinas rotativas. A vibração mecânica
é consequência direta das forças dinâmicas internas da máquina e, portanto,
carrega informação sobre o estado de componentes como rotores, mancais,
acoplamentos e rolamentos \cite{harris2002}. A análise espectral por
Transformada Rápida de Fourier (FFT) permite decompor o sinal de vibração e
identificar componentes de frequência cuja origem cinemática pode ser
associada a defeitos específicos \cite{randall2011,cooley1965}. Essa relação
entre componente espectral e fenômeno mecânico fundamenta o diagnóstico de
falhas em máquinas rotativas.

O avanço da chamada Indústria 4.0, contudo, trouxe um desafio paradoxal: a
digitalização e a sensorização massiva produzem um volume de dados brutos que,
se armazenado de forma indiscriminada, sobrecarrega as infraestruturas e
transforma um ativo analítico em passivo de armazenamento. Em sistemas de
turbomáquinas, frequências de interesse da ordem de 5~kHz exigem taxas de
amostragem de dezenas de milhares de amostras por segundo, multiplicadas por
centenas de pontos de medição. Esse cenário motivou o desenvolvimento de
estratégias de \emph{persistência seletiva}, em que a relevância técnica do
dado precede a sua gravação, e o dado bruto redundante é substituído por
representações compactas no domínio da frequência.

Paralelamente, o treinamento de analistas de vibração e de algoritmos de
inteligência artificial (IA) para diagnóstico depende da disponibilidade de
sinais de defeitos variados, em diferentes severidades e condições
operacionais. A aquisição desses sinais em ativos reais é lenta, cara e,
frequentemente, inviável: um defeito deve primeiro ocorrer, ser detectado e ser
documentado. A \emph{geração de sinais sintéticos} surge, então, como
alternativa para emular falhas de forma controlada e segura, permitindo testar
algoritmos e treinar analistas sem expor ativos críticos a riscos.

Este trabalho situa-se na interseção desses dois problemas. A partir de um
conjunto de diretrizes técnicas de engenharia de confiabilidade --- que
estabelecem critérios metrológicos de aquisição, um critério de persistência
seletiva denominado \(R^3\) e um catálogo de assinaturas espectrais de defeitos
--- foi especificado, projetado e implementado o \textbf{ARGUS}, um simulador de
defeitos em máquinas rotativas com geração de sinais vibracionais sintéticos,
análise espectral por FFT e persistência seletiva de leituras, acompanhado de
uma interface web para análise visual e registro de \emph{snapshots} de
defeito.

\section{Problema}

A análise do contexto evidencia dois problemas interligados que este trabalho
aborda:

\begin{enumerate}
  \item \textbf{Volume de dados brutos de vibração.} A aquisição contínua de
  sinais no domínio do tempo gera um volume massivo de dados. Em frequências
  de interesse de até 5~kHz, a taxa de amostragem mínima, segundo o critério de
  Nyquist, é da ordem de dezenas de milhares de amostras por segundo, o que
  torna o armazenamento indiscriminado inviável em escala industrial. É
  necessário reduzir o volume persistido sem perder a capacidade diagnóstica,
  conservando apenas a informação relevante --- os picos espectrais e as
  métricas de energia do sinal --- e descartando a redundância.
  \item \textbf{Falta de sinais de defeito para treinamento e pesquisa.}
  Analistas e algoritmos de IA precisam de exemplos de defeitos em diferentes
  estágios de severidade para aprendizado e validação. A obtenção desses
  exemplos em ativos reais é limitada por risco, indisponibilidade e custo,
  motivando a geração sintética de sinais cujas assinaturas espectrais sejam
  tecnicamente coerentes com a literatura de diagnóstico de vibração.
\end{enumerate}

O desafio de engenharia é, portanto, duplo: \emph{como gerar} sinais
vibracionais sintéticos cujas assinaturas espectrais representem defeitos
específicos de máquinas rotativas, e \emph{como persistir seletivamente} as
leituras de forma que apenas as transições de estado relevantes sejam
armazenadas, preservando a rastreabilidade para a análise de causa raiz (RCA).

\section{Justificativa}

A justificativa deste trabalho apoia-se em três pilares. O primeiro é
estratégico e econômico: a redução do volume de dados persistidos reduz o
custo de infraestrutura de armazenamento e direciona a análise para o que é
realmente crítico, aumentando o retorno sobre o investimento da manutenção
preditiva \cite{scheffer2004}. O segundo é metodológico: a capacidade de
emular falhas permite validar algoritmos de diagnóstico e treinar analistas sem
expor ativos reais, o que acelera o desenvolvimento de competências e de
modelos de IA. O terceiro é tecnológico: o trabalho materializa, em um sistema
funcional e testado, um conjunto de diretrizes de engenharia de confiabilidade
que, sem uma implementação de referência, permaneceriam apenas conceituais.

Além disso, o projeto foi conduzido seguindo uma metodologia de
especificação dirigida (Spec-Driven Development), o que gerou um acervo
documental --- especificação de requisitos, design, arquitetura, testes --- que
constitui evidência verificável do processo de engenharia adotado e base para a
reprodutibilidade dos resultados aqui apresentados.

\section{Objetivos}

\subsection{Objetivo geral}

Projetar e implementar um simulador de defeitos em máquina rotativa que gere
sinais vibracionais sintéticos com assinatura espectral parametrizável por tipo
de defeito e que aplique persistência seletiva de leituras --- por meio da
regra de ouro e do Paralelepípedo de Descarte no critério \(R^3\) --- com uma
interface web para análise espectral e registro de \emph{snapshots} de defeito
para fins de treinamento e pesquisa.

\subsection{Objetivos específicos}

\begin{enumerate}
  \item Modelar e implementar a geração de sinais sintéticos cujas componentes
  espectrais (ordens síncronas, harmônicos, sub-harmônicos, fracionários e
  frequências de falha de rolamento) reproduzam o catálogo de defeitos de
  máquinas rotativas adotado neste trabalho.
  \item Implementar o processamento espectral por FFT com extração de picos no
  espaço \(R^3\) (frequência, amplitude e fase), cálculo de RMS total, de ruído
  e de picos, valor DC e limiar configurável de detecção de picos.
  \item Implementar o motor de descarte (regra de ouro + Paralelepípedo de
  Descarte \(R^3\)) e a persistência seletiva com hierarquia
  Planta~$\rightarrow$~Área~$\rightarrow$~Máquina~$\rightarrow$~Ponto e
  armazenamento de leituras descartadas (\emph{trash}).
  \item Expor uma API REST e uma interface web com painel de sinal, gráfico de
  FFT com ordens normalizadas \(N\) e rotação em Hz, indicadores de tempo de
  processamento e taxa de descarte e funcionalidade de \emph{snapshot} de
  defeito.
  \item Validar o sistema por testes automatizados (unitários, de integração e
  de contrato) e por ensaio de ponta a ponta, verificando as assinaturas
  espectrais geradas e o fluxo de descarte.
\end{enumerate}

\section{Delimitação}

O trabalho delimita-se ao desenvolvimento de um \emph{simulador} de sinais, e
não de um sistema de aquisição real. As seguintes fronteiras são assumidas:

\begin{itemize}
  \item os sinais são gerados sinteticamente e não correspondem a medições em
  ativos reais; as amplitudes relativas dos perfis são parâmetros de síntese,
  conforme discutido no Capítulo~\ref{cap:fundamentacao};
  \item as frequências de falha de rolamento (BPFO, BPFI, BSF e FTF) são
  calculadas a partir de uma geometria padrão de rolamento;
  \item o ambiente de execução é local (Docker Compose), sem autenticação e sem
  ambiente de produção definidos, conforme decisões de projeto registradas no
  Capítulo~\ref{cap:desenvolvimento};
  \item o processamento é síncrono por requisição, sem filas ou lotes.
\end{itemize}

\section{Síntese da metodologia}

Do ponto de vista metodológico, o trabalho classifica-se como pesquisa
aplicada e desenvolvimento experimental, combinando engenharia de software,
processamento digital de sinais e engenharia de confiabilidade. O processo
adotado foi a especificação dirigida (Spec-Driven Development), no qual os
requisitos funcionais e não funcionais são formalizados antes da implementação,
o design é detalhado por tarefas pequenas e rastreáveis, e cada módulo é
validado por testes automatizados. A metodologia completa, incluindo materiais,
ambiente e critérios de aceitação, é descrita no Capítulo~\ref{cap:metodologia}.

\section{Organização do trabalho}

Além desta introdução, o trabalho está organizado da seguinte forma:

\begin{itemize}
  \item o \textbf{Capítulo~\ref{cap:fundamentacao}} apresenta a fundamentação
  teórica: vibração em máquinas rotativas, aquisição e processamento de sinais,
  FFT e análise espectral, catálogo de defeitos e o critério de persistência
  seletiva \(R^3\);
  \item o \textbf{Capítulo~\ref{cap:trabalhos}} discute trabalhos relacionados
  e posiciona a contribuição deste trabalho;
  \item o \textbf{Capítulo~\ref{cap:metodologia}} descreve materiais e métodos,
  incluindo o processo de engenharia e os critérios de validação;
  \item o \textbf{Capítulo~\ref{cap:desenvolvimento}} detalha a arquitetura e a
  implementação do sistema;
  \item o \textbf{Capítulo~\ref{cap:experimentos}} apresenta os experimentos e
  resultados;
  \item o \textbf{Capítulo~\ref{cap:discussao}} discute os resultados e as
  limitações;
  \item o \textbf{Capítulo~\ref{cap:conclusao}} conclui o trabalho e indica
  trabalhos futuros.
\end{itemize}

% =====================================================================
%  Capítulo 2 — Fundamentação Teórica
% =====================================================================
\chapter{Fundamentação Teórica}
\label{cap:fundamentacao}

Este capítulo apresenta os conceitos necessários à compreensão do problema e
da solução desenvolvidos neste trabalho: vibração mecânica em máquinas
rotativas, aquisição e processamento digital de sinais, análise espectral por
FFT, catálogo de defeitos e assinaturas espectrais e o critério de
persistência seletiva \(R^3\). Os conceitos são conectados às diretrizes
técnicas que serviram de especificação para o sistema ARGUS.

\section{Vibração mecânica em máquinas rotativas}

A vibração mecânica é a resposta dinâmica de uma estrutura às forças
alternadas que atuam sobre ela. Em máquinas rotativas, essas forças têm origem
em fenômenos como desbalanceamento residual do rotor, desalinhamento de
eixos, folgas mecânicas, contato rotor-estator e defeitos de mancais
\cite{harris2002,randall2011}. Como cada fenômeno excita a estrutura em
frequências características, o sinal de vibração medido em um ponto da máquina
contém informação diagnóstica sobre o estado dos seus componentes.

A grandeza fundamental para a interpretação espectral é a \emph{frequência de
rotação} do eixo. Dada a rotação em rotações por minuto (RPM), a frequência
fundamental, em hertz, é
\begin{equation}
\label{eq:fr}
  f_r = \frac{RPM}{60}.
\end{equation}
As componentes espectrais de um defeito síncrono aparecem em múltiplos
inteiros ou fracionários dessa frequência fundamental. A componente de ordem
\(k\) (denotada \(kX\)) possui frequência
\begin{equation}
\label{eq:fk}
  f_{kX} = k\, f_r .
\end{equation}
Por exemplo, a componente de primeira ordem (\(1X\)) ocorre exatamente na
frequência de rotação; a de segunda ordem (\(2X\)) no dobro, e assim por
diante. A representação por \emph{ordens} tem a vantagem de ser independente
da rotação absoluta: um defeito caracterizado por \(2X\) dominante mantém essa
característica em qualquer regime de rotação, desde que as frequências sejam
recalculadas pela Equação~\eqref{eq:fr}. Essa é a convenção adotada tanto na
literatura de diagnóstico quanto na especificação do ARGUS, na qual todas as
frequências são derivadas da rotação e nunca armazenadas como valores
absolutos fixos.

A norma ISO~10816 estabelece critérios para a avaliação da severidade da
vibração medida em partes não rotativas de máquinas, e a qualificação de
analistas de monitoramento de condição por vibração é tratada pela série
ISO~18436 \cite{iso10816,iso18436}. A convenção de unidades adotada depende da
grandeza medida e da faixa de frequência de interesse, conforme a
Tabela~\ref{tab:unidades}, que reproduz a diretriz metrológica que originou o
projeto.

\begin{table}[htbp]
\centering
\caption{Parâmetro, unidade de engenharia, convenção e aplicação típica na análise de vibração.}
\label{tab:unidades}
\begin{tabular}{@{}llll@{}}
\toprule
\textbf{Parâmetro} & \textbf{Unidade} & \textbf{Convenção} & \textbf{Aplicação típica} \\
\midrule
Deslocamento & mils / $\mu$m & Pk-Pk  & Baixas frequências; mancais de deslizamento; órbitas \\
Velocidade   & in/s / mm/s  & Peak   & Severidade global; falhas mecânicas (1X, 2X, 3X) \\
Aceleração   & g's / m/s²   & RMS    & Altas frequências; rolamentos; engrenamento; cavitação \\
\bottomrule
\end{tabular}
\caption*{\footnotesize Fonte: diretrizes técnicas do projeto (docs/descricao.txt).}
\end{table}

\section{Aquisição de sinais de vibração}

A qualidade do diagnóstico depende diretamente da qualidade da aquisição. Dois
aspectos são centrais: o critério de amostragem e o condicionamento do sinal.

\subsection{Critério de Nyquist e resolução espectral}

Para digitalizar um sinal analógico sem perda de informação, a taxa de
amostragem \(f_s\) deve ser superior ao dobro da maior frequência de interesse
\(F_{max}\) (critério de Nyquist),
\begin{equation}
\label{eq:nyquist}
  f_s > 2\, F_{max},
\end{equation}
sendo o uso de filtros \emph{anti-aliasing} de hardware mandatório para impedir
que energias acima de \(f_s/2\) se \emph{aliasem} como componentes de baixa
frequência, mascarando inclusive sinais subsíncronos \cite{oppenheim2010}.
Quando \(F_{max}\) é definido pela ordem harmônica mais alta do defeito de
interesse, o critério da Equação~\eqref{eq:nyquist} determina o limite mínimo
de amostragem. No ARGUS, essa validação é feita de forma automática antes de
qualquer processamento (comportamento \emph{fail-fast}), rejeitando
configurações que violem o critério.

A \emph{resolução espectral} \(\Delta f\) --- a menor separação entre duas
frequências distinguíveis no espectro --- é determinada pela taxa de amostragem
e pelo número de amostras \(N\) do registro:
\begin{equation}
\label{eq:df}
  \Delta f = \frac{f_s}{N}.
\end{equation}
Uma resolução adequada é necessária para separar picos adjacentes e para
identificar bandas laterais (\emph{sidebands}), indicadores importantes de
severidade em falhas de rolamento e de modulação \cite{randall2011}.

\subsection{Janelamento}

Como a FFT pressupõe que o sinal seja periódico no registro analisado, o
truncamento de um sinal não periódico introduz vazamento espectral
(\emph{leakage}): a energia de uma componente se espalha para frequências
vizinhas. O janelamento pondera as amostras para mitigar esse efeito. A
Tabela~\ref{tab:janelas} resume os critérios de aplicação adotados como
referência metrológica.

\begin{table}[htbp]
\centering
\caption{Tipos de janela e critérios de aplicação no processamento espectral.}
\label{tab:janelas}
\begin{tabular}{@{}p{3cm}p{5.5cm}p{4.6cm}@{}}
\toprule
\textbf{Janela} & \textbf{Critério de aplicação} & \textbf{Impacto na precisão} \\
\midrule
Hanning   & Padrão para sinais contínuos em regime permanente & Equilíbrio entre seletividade e proteção contra dispersão \\
Flat Top  & Obrigatória para testes de aceitação e calibração & Prioriza a precisão absoluta da amplitude do pico \\
Rectangular & Restrita a sinais transientes que iniciam e terminam no registro & Eficaz apenas se o registro contém o transiente completo \\
\bottomrule
\end{tabular}
\caption*{\footnotesize Fonte: diretrizes técnicas do projeto (docs/descricao.txt).}
\end{table}

\section{Transformada de Fourier e análise espectral}

A análise espectral converte o sinal do domínio do tempo para o domínio da
frequência, revelando as componentes que compõem a vibração. Para um sinal
discreto \(x[n]\) com \(N\) amostras, a Transformada Discreta de Fourier (DFT)
é definida por
\begin{equation}
\label{eq:dft}
  X[k] = \sum_{n=0}^{N-1} x[n]\, e^{-j 2\pi k n / N}, \qquad k = 0,\ldots,N-1,
\end{equation}
em que \(X[k]\) é a componente espectral complexa do índice \(k\). A magnitude
\(|X[k]|\) indica a amplitude da componente na frequência \(k\,\Delta f\), e o
argumento \(\angle X[k]\) indica a sua fase. A Transformada Rápida de Fourier
(FFT) é um algoritmo eficiente para o cálculo da DFT, com complexidade
\(O(N \log N)\), viabilizando a análise espectral em tempo real ou quase real
\cite{cooley1965,oppenheim2010}.

A amplitude espectral de um pico corresponde, no domínio do tempo, à amplitude
de uma senoide de frequência igual à do pico. A fase, por sua vez, codifica a
posição angular da componente em relação a uma referência de tempo --- o
\emph{keyphasor} ou marcador de rotação --- e é essencial para distinguir
fenômenos como desbalanceamento estático e de acoplamento, que apresentam
relações de fase distintas entre mancais \cite{randall2011,harris2002}.

\section{O espaço \(R^3\) e as métricas do sinal}
\label{sec:r3}

A diretriz de persistência seletiva que fundamenta este trabalho propõe
representar cada componente espectral relevante como um ponto no espaço
tridimensional \(R^3\):
\begin{itemize}
  \item \textbf{frequência} --- identifica a origem cinemática do fenômeno
  (ordem \(kX\) em relação à rotação);
  \item \textbf{amplitude} --- quantifica a severidade do esforço mecânico;
  \item \textbf{fase} --- provê a correlação temporal para diagnósticos como
  desbalanceamento e análise orbital.
\end{itemize}

Além dos picos, o sistema calcula métricas de energia do sinal. O valor eficaz
(RMS) total de um sinal com \(N\) amostras é
\begin{equation}
\label{eq:rms}
  x_{RMS} = \sqrt{\frac{1}{N}\sum_{n=0}^{N-1} x[n]^2}.
\end{equation}
O RMS é a métrica padrão para a severidade global em aceleração e é a base de
critérios normativos como a ISO~10816 \cite{iso10816}. No contexto deste
trabalho, distinguem-se o RMS total, o RMS das componentes de pico (calculado a
partir das amplitudes dos picos \(R^3\)) e o RMS do ruído de fundo (obtido por
diferença de energia entre o total e os picos). O \emph{valor DC} é a média do
sinal e corresponde ao \emph{gap}/bias de sensores de proximidade (medição de
deslocamento), sendo persistido para permitir reconstrução de tendências.

\section{Catálogo de defeitos e assinaturas espectrais}

A análise de vibração relaciona padrões espectrais a defeitos mecânicos
específicos. O projeto adota um catálogo com treze tipos de defeito
(Tabela~\ref{tab:catalogo}), com assinaturas derivadas da literatura de
diagnóstico \cite{randall2011,scheffer2004} e formalizadas em um documento de
especificação de assinaturas utilizado como entrada para a geração sintética de
sinais.

\begin{table}[htbp]
\centering
\caption{Catálogo de tipos de defeito e assinatura espectral dominante.}
\label{tab:catalogo}
\begin{tabular}{@{}p{4.4cm}p{8.6cm}@{}}
\toprule
\textbf{Tipo} & \textbf{Assinatura dominante} \\
\midrule
sem\_defeito & 1X residual baixo, sem sequência forte de harmônicos \\
desbalanceamento & 1X dominante \\
desalinhamento\_angular & 2X + 1X, com 3X e fracionários 1,5X/2,5X \\
desalinhamento\_paralelo & 2X dominante + harmônicos pares 4X/6X/8X \\
rocamento & 1X + série de harmônicos; sub-harmônicos (0,5X, 1/3X, 2/3X) quando severo \\
mancal\_frouxo & família de harmônicos 1X--10X + fracionários \\
acoplamento\_defeituoso & 1X/2X + 3X/4X \\
oil\_whirl & subsíncrona 0,39X--0,48X (razão amostrada) + 1X \\
whirl\_atrito & subsíncrona não universal (faixa paramétrica, hipótese) \\
rolamento\_bpfo & BPFO + harmônicos $\pm$ sidebands 1X \\
rolamento\_bpfi & BPFI + harmônicos $\pm$ sidebands 1X \\
rolamento\_bsf & BSF + harmônicos \\
rolamento\_ftf & FTF + harmônicos \\
\bottomrule
\end{tabular}
\caption*{\footnotesize Fonte: SPECIFICATION.md do projeto; docs/assinaturas\_fft\_falhas\_maquinas\_rotativas\_IA.md.}
\end{table}

\subsection{Defeitos síncronos e harmônicos}

Forças centrífugas de \emph{desbalanceamento} crescem com o quadrado da
rotação e manifestam-se pela dominância de \(1X\) na direção radial; conforme o
tipo (estático, de acoplamento ou dinâmico), a fase entre mancais distingue o
quadro \cite{harris2002,randall2011}. O \emph{desalinhamento} angular produz
predominância de \(1X\) e \(2X\) na direção axial, com defasagem de \(180^\circ\)
através do acoplamento; o desalinhamento paralelo manifesta-se pela dominância
de \(2X\) na direção radial, com harmônicos pares superiores em casos severos.
\emph{Folgas mecânicas} (looseness) do tipo estrutural elevam \(1X\) vertical,
enquanto folgas móveis ou de mancal geram uma ``floresta'' de harmônicos
síncronos de \(1X\) até \(10X\). O \emph{roçamento} de rotor (rub) produz um
espectro rico em harmônicos e, distintamente, sub-harmônicos (\(1/2X\), \(1/3X\)),
com elevação do piso de ruído causada pelo truncamento da forma de onda
\cite{scheffer2004}.

\subsection{Instabilidades de filme de óleo}

Em mancais de deslizamento, a instabilidade de \emph{oil whirl} manifesta-se
por uma componente subsíncrona na faixa de aproximadamente \(0{,}40X\) a
\(0{,}49X\) da rotação, com órbita circular. Quando a frequência de whirl
coincide com uma frequência natural do sistema e ``trava'' nela, ocorre o
\emph{oil whip}, uma ressonância autoexcitada de caráter catastrófico
\cite{randall2011}. No catálogo do ARGUS, o \emph{oil whirl} é sintetizado com
razão amostrada uniformemente em \(0{,}39X\)--\(0{,}48X\) e o \emph{whirl por
atrito} como faixa paramétrica hipotética, sem razão universal.

\subsection{Falhas de rolamento}

Defeitos de mancais de rolamento geram frequências não síncronas características
que dependem da geometria do rolamento e da rotação. As frequências de falha
são, para um rolamento com \(n\) elementos rolantes, diâmetro do elemento
\(B_d\), diâmetro primitivo \(P_d\) e ângulo de contato \(\phi\)
\cite{randall2011,harris2002}:
\begin{align}
\label{eq:bpfo}
  BPFO &= \frac{n}{2}\, f_r \left(1 - \frac{B_d}{P_d}\cos\phi\right),\\[2pt]
\label{eq:bpfi}
  BPFI &= \frac{n}{2}\, f_r \left(1 + \frac{B_d}{P_d}\cos\phi\right),\\[2pt]
\label{eq:bsf}
  BSF  &= \frac{P_d}{2 B_d}\, f_r \left[1 - \left(\frac{B_d}{P_d}\cos\phi\right)^2\right],\\[2pt]
\label{eq:ftf}
  FTF  &= \frac{1}{2}\, f_r \left(1 - \frac{B_d}{P_d}\cos\phi\right),
\end{align}
em que BPFO corresponde a defeito na pista externa, BPFI na pista interna, BSF
nos elementos rolantes e FTF na gaiola. O surgimento de \emph{sidebands} em
torno dessas frequências indica modulação (tipicamente em \(1X\)) e severidade
avançada \cite{randall2011}. O ARGUS calcula essas ordens a partir de uma
geometria padrão (\(n=9\), \(B_d/P_d=0{,}2\), \(\phi=0\)), reproduzindo as
Equações~\eqref{eq:bpfo}--\eqref{eq:ftf} como múltiplos da frequência de
rotação; a geometria real do rolamento é uma pendência para uso em campo.

\subsection{Observação sobre amplitudes sintéticas}

As amplitudes relativas dos perfis adotados são \emph{parâmetros de síntese}
para geração de sinais e conjuntos de dados, e não valores universais de
severidade: a amplitude real depende de massa, rigidez, amortecimento,
posição do sensor, carga, ressonâncias e montagem. Essa distinção é essencial
para o uso academicamente honesto do simulador como ferramenta de treinamento,
e será retomada na discussão (Capítulo~\ref{cap:discussao}).

\section{Persistência seletiva e o critério \(R^3\)}

O segundo pilar do problema é o armazenamento de dados. A monitoração contínua
de turbomáquinas, com frequências de interesse de até 5~kHz, exige taxas de
amostragem entre 12,5~kHz e 50~kHz; multiplicadas por centenas de pontos de
medição, essas taxas geram um volume que inviabiliza o armazenamento
indiscriminado. A estratégia adotada --- e materializada no ARGUS --- é a
\emph{persistência seletiva}: migrar do domínio do tempo para o domínio da
frequência, conservando os vetores de picos \(R^3\) e as métricas de energia, e
descartar leituras redundantes cujos picos permaneçam dentro de uma zona de
tolerância em relação à última leitura efetivamente persistida.

A decisão de persistir ou descartar segue uma hierarquia lógica que evita o
\emph{drift} de dados (perda de pequenas variações acumuladas que seriam
ignoradas se a comparação fosse feita apenas contra a leitura anterior, que pode
ter sido descartada):

\begin{enumerate}
  \item \textbf{Variação de rotação} --- detecta mudança de regime operacional
  (rampas, carga). Pela \emph{regra de ouro}, a leitura é persistida se
  \begin{equation}
  \label{eq:regra-ouro}
    \left|RPM_{atual} - RPM_{armazenada}\right| > \Delta_{rot}
    \;\wedge\; (\text{número de picos} > 0).
  \end{equation}
  \item \textbf{Número de picos} --- identifica o surgimento de novas
  componentes espectrais.
  \item \textbf{Verificação \(R^3\)} --- aplica o \emph{Paralelepípedo de
  Descarte}: a leitura é descartada somente se, para cada pico atual, existir um
  pico de referência (o mais próximo em frequência, na última leitura
  persistida) tal que
  \begin{equation}
  \label{eq:descarte}
    |\Delta f| \le \delta_f \;\wedge\; |\Delta A| \le \delta_A
    \;\wedge\; |\Delta \phi| \le \delta_\phi,
  \end{equation}
  em que \(\delta_f\), \(\delta_A\) e \(\delta_\phi\) são as tolerâncias em
  frequência, amplitude e fase. Se ao menos um pico estiver fora dessa zona, a
  leitura é persistida e passa a ser a nova referência.
\end{enumerate}

A primeira leitura de um ponto é sempre persistida (condição de inicialização),
e a comparação é sempre feita contra a última leitura \emph{efetivamente
persistida}, nunca contra a última leitura avaliada/descartada. As leituras
descartadas são registradas em um armazenamento separado (\emph{trash}),
permitindo auditar a taxa de descarte e validar o algoritmo. Esse critério
viabiliza a redução do volume de dados sem perda da capacidade diagnóstica,
direcionando o esforço analítico para as transições de estado relevantes do
ativo \cite{scheffer2004}.

\section{Considerações normativas}

O trabalho observa, na sua apresentação, as normas ABNT de documentação
(NBR~14724 e NBR~6023) \cite{abnt14724,abnt6023} e, na fundamentação técnica,
as normas ISO de avaliação de vibração e qualificação de pessoal
\cite{iso10816,iso18436}. As edições vigentes dessas normas na data da defesa
devem ser verificadas, conforme registrado nas pendências do projeto.

% =====================================================================
%  Capítulo 3 — Trabalhos relacionados
% =====================================================================
\chapter{Trabalhos Relacionados}
\label{cap:trabalhos}

Este capítulo situa o trabalho em relação à literatura e a sistemas correlatos.
A revisão aqui apresentada é delimitada (não sistemática) e foca nas três
dimensões do problema: monitoramento de condição por vibração, geração de
sinais sintéticos e conjuntos de dados para treinamento, e estratégias de
armazenamento seletivo de dados de vibração.

\section{Monitoramento de condição e análise de vibração}

O monitoramento de condição por vibração é um campo maduro, consolidado em
normas e em sistemas comerciais. Randall \cite{randall2011} apresenta o estado
da arte das técnicas de monitoramento baseado em vibração para aplicações
industriais, aeroespaciais e automotivas, incluindo a análise de envelope para
falhas incipientes de rolamento e a interpretação de bandas laterais. Scheffer
e Girdhar \cite{scheffer2004} oferecem uma abordagem prática da análise de
vibração e da manutenção preditiva, cobrindo o fluxo completo desde a
aquisição até o diagnóstico e a recomendação de ação. O manual de Harris e
Piersol \cite{harris2002} é referência clássica para as bases teóricas da
dinâmica e da instrumentação de vibração.

Sistemas comerciais de monitoramento (provenientes de fornecedores de
instrumentação e de plataformas de manutenção preditiva) implementam
funcionalidades semelhantes às deste trabalho --- aquisição, FFT, alarmes e
tendências ---, porém tipicamente como produtos fechados, acoplados a hardware
específico de aquisição. O diferencial do ARGUS não está na análise espectral em
si, mas na combinação de \emph{geração sintética de defeitos} e
\emph{persistência seletiva} em um sistema aberto, extensível e validado por
testes, orientado ao treinamento e à pesquisa.

\section{Bases de dados e sinais sintéticos para treinamento}

Para treinamento de analistas e de algoritmos de aprendizagem de máquina, a
comunidade utiliza tanto bases de dados medidas em bancadas de ensaio quanto
sinais sintéticos. As bases medidas, ainda que reais, apresentam limitações
importantes: cobrem um conjunto restrito de máquinas, condições e defeitos;
exigem a reprodução física da falha (com custo e risco); e a documentação dos
estágios de severidade nem sempre é completa. Sinais sintéticos, por sua vez,
permitem gerar grandes volumes rotulados com controle de severidade, de ruído e
de condição operacional, a custo desprezível e sem risco --- exatamente o
propósito do gerador desenvolvido neste trabalho. A limitação inerente é que as
amplitudes sintéticas são parâmetros de síntese e não refletem a dinâmica real
de uma máquina específica (massas, rigidez, amortecimento, ressonâncias), como
discutido no Capítulo~\ref{cap:fundamentacao} e retomado no
Capítulo~\ref{cap:discussao}.

A especificação de assinaturas adotada no projeto segue a literatura de
diagnóstico para definir quais componentes espectrais (ordens e frequências de
falha) devem estar presentes em cada defeito, o que confere coerência física aos
sinais gerados, ainda que as amplitudes sejam relativas.

\section{Persistência seletiva de dados de vibração}

O problema da ``massa de dados'' no monitoramento de condição é reconhecido na
indústria, e a literatura de manutenção preditiva aponta a necessidade de
transformar o dado bruto em informação antes do armazenamento
\cite{scheffer2004}. A abordagem do ARGUS --- representar a leitura por vetores
de picos \(R^3\) e métricas de energia, e persistir apenas transições de estado
segundo um Paralelepípedo de Descarte --- segue essa diretriz e a materializa
em um motor de descarte testável. A comparação contra a última leitura
\emph{efetivamente persistida} (e não contra a última leitura avaliada) é um
detalhe de projeto que evita o acúmulo de \emph{drift}, aspecto que distingue a
implementação de um descarte ingênuo por diferença da leitura anterior.

\section{Lacuna e contribuição deste trabalho}

A partir da revisão delimitada, identifica-se a seguinte lacuna: ferramentas
abertas que integrem, de forma validada e testável, (i) a geração sintética de
sinais de vibração com assinaturas espectrais de um catálogo amplo de defeitos,
(ii) a persistência seletiva segundo o critério \(R^3\) e (iii) uma interface
web para análise e registro de \emph{snapshots} de defeito. O ARGUS propõe-se a
preencher essa lacuna como uma plataforma de treinamento e pesquisa, apoiada em
documentação de requisitos, design e testes que garantem reprodutibilidade.

Uma revisão bibliográfica sistemática e o aprofundamento quantitativo das
comparações com bases de dados públicas são registrados como pendências do
trabalho (ver `MONOGRAFIA_PENDENCIAS.md`), devendo ser conduzidos com o
orientador antes da versão final.

% =====================================================================
%  Capítulo 4 — Materiais e Métodos
% =====================================================================
\chapter{Materiais e Métodos}
\label{cap:metodologia}

Este capítulo descreve a natureza da pesquisa, o processo de engenharia
adotado, os materiais e o ambiente utilizados e os critérios de validação,
permitindo a reprodução ou a avaliação crítica do trabalho. Procura-se
distinguir claramente o que foi \emph{planejado} (nos artefatos de
especificação e design) do que foi \emph{executado} (código e resultados).

\section{Natureza e abordagem da pesquisa}

O trabalho caracteriza-se como \textbf{pesquisa aplicada} com componente de
\textbf{desenvolvimento experimental}: parte de um problema de engenharia real
(volume de dados e necessidade de sinais de defeito no monitoramento de
condição) e entrega um artefato funcional --- um simulador --- com validação
por testes. A abordagem é essencialmente tecnológica e combina três áreas:
engenharia de software, processamento digital de sinais e engenharia de
confiabilidade (manutenção preditiva).

\section{Processo de engenharia: especificação dirigida}

O desenvolvimento seguiu o processo de \emph{Spec-Driven Development} (SDD),
no qual o código é consequência de uma especificação rastreável, e não o
contrário. O fluxo executado foi:

\begin{enumerate}
  \item \textbf{Auditoria do contexto}: análise das diretrizes técnicas de
  origem (critérios metrológicos, critério \(R^3\), catálogo de defeitos),
  disponíveis como documentos de entrada do projeto.
  \item \textbf{Especificação}: formalização dos requisitos funcionais
  (FR-001 a FR-019), requisitos não funcionais (NFR-001 a NFR-007) e critérios
  de aceitação (CA-001 a CA-014) em \texttt{SPECIFICATION.md}, com matriz de
  rastreabilidade.
  \item \textbf{Design}: arquitetura, contratos de API, modelo de dados e
  detalhamento por tarefas pequenas (T-001 a T-013) em \texttt{.specs/}.
  \item \textbf{Implementação}: execução das tarefas com gate mínimo de
  lint, type-check e testes por módulo.
  \item \textbf{Validação}: testes unitários, de integração e de contrato, e
  ensaio de ponta a ponta no navegador.
\end{enumerate}

Cada etapa deixou evidência verificável (especificação, design, tarefas,
testes e registros de execução), o que é utilizado neste texto como base das
afirmações sobre o sistema.

\section{Materiais e ambiente}

O sistema foi desenvolvido e validado no ambiente descrito na
Tabela~\ref{tab:ferramentas}. As versões informadas correspondem às
declaradas nos artefatos do projeto e verificadas durante a execução.

\begin{table}[htbp]
\centering
\caption{Materiais, ferramentas e versões utilizados no desenvolvimento e na validação.}
\label{tab:ferramentas}
\begin{tabular}{@{}lll@{}}
\toprule
\textbf{Camada} & \textbf{Ferramenta} & \textbf{Versão/nota} \\
\midrule
Back-end & Python & 3.12 \\
         & FastAPI & $\ge$ 0.115 (Pydantic v2) \\
         & NumPy / SciPy & $\ge$ 1.26 / $\ge$ 1.13 \\
         & SQLAlchemy + asyncpg + Alembic & 2.x / $\ge$ 0.29 / $\ge$ 1.13 \\
         & ruff / mypy / pytest & lint, format, type-check, testes \\
Front-end & React & 19.2.x instalada \\
          & Vite / TypeScript & 8.2.x / 6.0.x \\
          & Recharts & 3.10.x \\
          & oxlint / Vitest / Testing Library & lint e testes \\
Banco & PostgreSQL & 16 (contêiner Docker) \\
Infra & Docker / Docker Compose & v2 \\
Execução & Navegador & ensaio E2E em localhost:5173 \\
\bottomrule
\end{tabular}
\par\vspace{2pt}{\footnotesize Fonte: \texttt{frontend/package.json}, \texttt{backend/requirements*.txt} e \texttt{docker-compose.yml} do projeto.}
\end{table}

\section{Arquitetura de teste em camadas}

A estratégia de validação seguiu uma pirâmide de testes, com camadas de custo
crescente e escopo decrescente:

\begin{enumerate}
  \item \textbf{Testes unitários} dos módulos puros de domínio
  (\texttt{signal\_generator}, \texttt{fft\_processor}, \texttt{discard\_engine}), sem acesso a banco
  de dados ou rede. O motor de descarte recebe a última leitura persistida como
  parâmetro injetado, o que permite cobrir todas as combinações da regra de ouro
  e do Paralelepípedo de Descarte sem infraestrutura.
  \item \textbf{Testes de integração} da camada de persistência, executados
  contra um PostgreSQL real provisório (contêiner via \texttt{testcontainers}),
  validando migração, transações e o ponteiro de última leitura persistida.
  \item \textbf{Testes de contrato} da API, usando cliente HTTP assíncrono
  contra a aplicação FastAPI, verificando os contratos de \texttt{POST /simulacoes},
  \texttt{POST /snapshots} e \texttt{GET /health}, incluindo validações de entrada e a
  violação de Nyquist.
  \item \textbf{Testes de front-end} (componentes e validação de formulário)
  com Vitest e Testing Library.
  \item \textbf{Ensaio de ponta a ponta} no navegador, contra o ambiente
  Docker Compose completo, verificando o fluxo real de simulação e a renderização
  dos elementos visuais.
\end{enumerate}

Os critérios de aceitação funcionais (CA-001 a CA-014) foram usados como
fonte dos casos de teste, ligando cada teste ao requisito que implementa.

\section{Métricas e critérios de validação}

As métricas observadas na validação foram:

\begin{itemize}
  \item \textbf{correção espectral}: presença e ordem das componentes
  esperadas para cada defeito do catálogo, verificadas no espectro retornado
  pela API;
  \item \textbf{correção do descarte}: decisão (persistir/descartar) de acordo
  com a regra de ouro e o Paralelepípedo \(R^3\), conforme os casos de teste;
  \item \textbf{tempo de processamento} por simulação (indicador exposto pela
  API, FR-017);
  \item \textbf{taxa de descarte acumulada} (leituras descartadas divididas
  pelas leituras avaliadas, NFR-004);
  \item \textbf{gates de qualidade}: lint e type-check limpos e suíte de
  testes verde nas duas camadas.
\end{itemize}

\section{Protocolo experimental executado}

O protocolo efetivamente executado foi o seguinte (distinguindo-se o planejado
do executado):

\begin{enumerate}
  \item subir o ambiente completo com \texttt{start.sh} (Docker Compose: banco,
  migrações Alembic, back-end e front-end);
  \item semear um ponto de medição de demonstração com identificador fixo;
  \item executar a suíte de testes do back-end (lint + type-check + pytest) e
  do front-end (type-check + lint + vitest);
  \item verificar, via chamadas reais à API, as assinaturas espectrais de
  defeitos representativos (ex.: desbalanceamento com \(1X\) dominante,
  desalinhamento paralelo com \(2X\)/\(4X\)/\(6X\), \emph{oil whirl} com
  razão subsíncrona e rolamento com BPFI e \emph{sidebands});
  \item executar o ensaio de ponta a ponta no navegador, simulando um defeito
  e verificando os elementos visuais do painel;
  \item registrar as evidências (contagens de testes, capturas de tela) para
  este texto.
\end{enumerate}

O item 2 (seed de ponto) foi executado por script de inicialização, uma vez que
o schema não expõe endpoint de criação de hierarquia no MVP; essa limitação é
discutida no Capítulo~\ref{cap:discussao}.

\section{Limitações metodológicas}

As principais limitações metodológicas são: (i) a validação é local e não há
ambiente de produção definido (NFR-005); (ii) a avaliação de desempenho
baseou-se em testes funcionais e observação direta, sem bateria estatística
formal de medições; e (iii) não foi conduzida revisão sistemática da literatura,
apenas revisão delimitada. Essas limitações estão registradas nas pendências do
trabalho e devem ser tratadas com o orientador.

% =====================================================================
%  Capítulo 5 — Desenvolvimento (arquitetura e implementação)
% =====================================================================
\chapter{Desenvolvimento}
\label{cap:desenvolvimento}

Este capítulo descreve a arquitetura e a implementação do ARGUS, relacionando
cada módulo aos requisitos que implementa. A descrição é fiel ao código
efetivamente produzido e testado, cujos artefatos principais residem em
`backend/` e `frontend/`.

\section{Visão geral da arquitetura}

O sistema é organizado em três camadas: \textbf{front-end} (aplica\c{c}\~ao web),
\textbf{back-end} (API e l\'ogica de dom\'inio) e \textbf{banco de dados}
(PostgreSQL). A Figura~\ref{fig:arquitetura} apresenta a arquitetura de
m\'odulos e os fluxos principais.

\begin{figure}[htbp]
\centering
% Diagrama: Arquitetura do ARGUS (visão de contêineres/módulos)
% Uso: % Diagrama: Arquitetura do ARGUS (visão de contêineres/módulos)
% Uso: \input{diagrams/arquitetura} dentro de um ambiente figure.
\begin{tikzpicture}[
  font=\small,
  fe/.style={rectangle, draw=black!60, fill=blue!6, rounded corners=2pt, align=center, inner sep=5pt, minimum width=3.8cm, text width=3.6cm},
  be/.style={rectangle, draw=black!60, fill=green!6, rounded corners=2pt, align=center, inner sep=5pt, minimum width=3.8cm, text width=3.6cm},
  db/.style={rectangle, draw=black!60, fill=orange!8, rounded corners=2pt, align=center, inner sep=5pt, minimum width=3.8cm, text width=3.6cm},
  lab/.style={font=\scriptsize},
  arrow/.style={-{Stealth[length=2.4mm]}, semithick}
]
  % Coluna front-end
  \node[fe] (ui) at (0,2.4) {Painel de simula\c{c}\~ao};
  \node[fe] (fftchart) at (0,0) {Gr\'afico de FFT + limiar\\ e ordens $N$};
  \node[fe] (check) at (0,-2.4) {Checklist de valida\c{c}\~ao};

  % Coluna back-end
  \node[be] (api) at (7,2.4) {API REST\\ \scriptsize POST /simulacoes, /snapshots};
  \node[be] (gen) at (7,0) {Gera\c{c}\~ao de sinal};
  \node[be] (fftm) at (7,-2.4) {FFT e extra\c{c}\~ao $R^3$};
  \node[be] (disc) at (7,-4.8) {Motor de descarte\\ \scriptsize regra de ouro + Paralelep\'ipedo};
  \node[be] (pers) at (7,-7.2) {Persist\^encia};

  % Coluna banco
  \node[db] (hier) at (14.5,3.6) {Planta/Área/\\M\'aquina/Ponto};
  \node[db] (snap) at (14.5,1.2) {snapshots\_defeito};
  \node[db] (leit) at (14.5,-1.2) {leituras\_persistidas};
  \node[db] (trash) at (14.5,-3.6) {leituras\_trash};

  % Fluxo front -> back
  \draw[arrow] (ui.east) -- node[lab,above]{POST /simulacoes} (api.west);
  \draw[arrow] (check.east) -- (api.west);
  \draw[arrow] (api.west) to[out=180,in=0] node[lab,above,near end]{resposta (sinal+FFT+RMS)} (fftchart.east);

  % Fluxo interno do back-end
  \draw[arrow] (api) -- (gen);
  \draw[arrow] (gen) -- (fftm);
  \draw[arrow] (fftm) -- (disc);
  \draw[arrow] (disc) -- (pers);

  % Fluxo back -> banco
  \draw[arrow] (api.east) to[out=0,in=180] (snap.west);
  \draw[arrow] (pers.east) -- node[lab,above,pos=0.4]{aprovada} (leit.west);
  \draw[arrow] (pers.east) -- node[lab,below,pos=0.3]{descartada} (trash.west);
  \draw[arrow] (pers.east) to[out=0,in=180] node[lab,above,pos=0.6]{l\^e/atualiza} (hier.west);
\end{tikzpicture}
 dentro de um ambiente figure.
\begin{tikzpicture}[
  font=\small,
  fe/.style={rectangle, draw=black!60, fill=blue!6, rounded corners=2pt, align=center, inner sep=5pt, minimum width=3.8cm, text width=3.6cm},
  be/.style={rectangle, draw=black!60, fill=green!6, rounded corners=2pt, align=center, inner sep=5pt, minimum width=3.8cm, text width=3.6cm},
  db/.style={rectangle, draw=black!60, fill=orange!8, rounded corners=2pt, align=center, inner sep=5pt, minimum width=3.8cm, text width=3.6cm},
  lab/.style={font=\scriptsize},
  arrow/.style={-{Stealth[length=2.4mm]}, semithick}
]
  % Coluna front-end
  \node[fe] (ui) at (0,2.4) {Painel de simula\c{c}\~ao};
  \node[fe] (fftchart) at (0,0) {Gr\'afico de FFT + limiar\\ e ordens $N$};
  \node[fe] (check) at (0,-2.4) {Checklist de valida\c{c}\~ao};

  % Coluna back-end
  \node[be] (api) at (7,2.4) {API REST\\ \scriptsize POST /simulacoes, /snapshots};
  \node[be] (gen) at (7,0) {Gera\c{c}\~ao de sinal};
  \node[be] (fftm) at (7,-2.4) {FFT e extra\c{c}\~ao $R^3$};
  \node[be] (disc) at (7,-4.8) {Motor de descarte\\ \scriptsize regra de ouro + Paralelep\'ipedo};
  \node[be] (pers) at (7,-7.2) {Persist\^encia};

  % Coluna banco
  \node[db] (hier) at (14.5,3.6) {Planta/Área/\\M\'aquina/Ponto};
  \node[db] (snap) at (14.5,1.2) {snapshots\_defeito};
  \node[db] (leit) at (14.5,-1.2) {leituras\_persistidas};
  \node[db] (trash) at (14.5,-3.6) {leituras\_trash};

  % Fluxo front -> back
  \draw[arrow] (ui.east) -- node[lab,above]{POST /simulacoes} (api.west);
  \draw[arrow] (check.east) -- (api.west);
  \draw[arrow] (api.west) to[out=180,in=0] node[lab,above,near end]{resposta (sinal+FFT+RMS)} (fftchart.east);

  % Fluxo interno do back-end
  \draw[arrow] (api) -- (gen);
  \draw[arrow] (gen) -- (fftm);
  \draw[arrow] (fftm) -- (disc);
  \draw[arrow] (disc) -- (pers);

  % Fluxo back -> banco
  \draw[arrow] (api.east) to[out=0,in=180] (snap.west);
  \draw[arrow] (pers.east) -- node[lab,above,pos=0.4]{aprovada} (leit.west);
  \draw[arrow] (pers.east) -- node[lab,below,pos=0.3]{descartada} (trash.west);
  \draw[arrow] (pers.east) to[out=0,in=180] node[lab,above,pos=0.6]{l\^e/atualiza} (hier.west);
\end{tikzpicture}

\caption{Arquitetura do ARGUS: front-end (React), back-end (FastAPI) e banco de dados (PostgreSQL), com fluxos de simulação, descarte e persistência.}
\label{fig:arquitetura}
\caption*{\footnotesize Fonte: elaboração própria, com base na arquitetura real do projeto.}
\end{figure}

O fluxo nominal de uma simulação é: (i) a interface envia os parâmetros a
`POST /simulacoes`; (ii) o back-end valida a entrada (incluindo o critério de
Nyquist), gera o sinal, calcula a FFT e extrai os picos \(R^3\) e as métricas;
(iii) o motor de descarte decide entre persistir e descartar com base na
última leitura persistida do ponto; (iv) a resposta retorna sinal decimado,
picos, métricas, decisão e indicadores. A Figura~\ref{fig:fluxo-descarte}
detalha o fluxo de decisão de descarte e a Figura~\ref{fig:estados} a máquina
de estados por leitura.

\begin{figure}[htbp]
\centering
% Diagrama: Fluxo de decisão do motor de descarte (regra de ouro + R^3)
% Uso: % Diagrama: Fluxo de decisão do motor de descarte (regra de ouro + R^3)
% Uso: \input{diagrams/fluxo-descarte} dentro de um ambiente figure.
\begin{tikzpicture}[
  font=\small,
  dec/.style={diamond, draw=black!60, fill=gray!8, align=center, inner sep=1pt, aspect=3, text width=2.2cm},
  act/.style={rectangle, draw=black!60, fill=green!6, rounded corners=2pt, align=center, inner sep=5pt, minimum width=4.0cm, text width=3.8cm},
  term/.style={rectangle, draw=black!60, fill=orange!12, rounded corners=2pt, align=center, inner sep=5pt, minimum width=4.2cm, text width=4.0cm},
  arrow/.style={-{Stealth[length=2.4mm]}, semithick},
  lab/.style={font=\scriptsize}
]
  \node[act] (inicio) at (0,3.0) {Leitura avaliada (rota\c{c}\~ao + picos $R^3$)};
  \node[dec] (primeira) at (0,0.6) {Primeira\\ leitura do\\ ponto?};
  \node[dec] (ouro) at (6.0,0.6) {$|\Delta RPM| > \Delta_{rot}$\\ e n\'umero de\\ picos $>0$?};
  \node[dec] (r3) at (12.0,0.6) {Todos os picos\\ dentro do\\ Paralelep\'ipedo?};
  \node[term] (persiste) at (6.0,-2.8) {PERSISTIR\\ \scriptsize leituras\_persistidas (nova refer\^encia)};
  \node[term] (descarta) at (12.0,-2.8) {DESCARTAR\\ \scriptsize registrar em leituras\_trash};

  \draw[arrow] (inicio) -- (primeira);

  \draw[arrow] (primeira.east) -- node[lab,above]{n\~ao} (ouro.west);
  \draw[arrow] (primeira.south) -- ++(0,-1.6) -| node[lab,above,pos=0.25]{sim} (persiste.west);

  \draw[arrow] (ouro.east) -- node[lab,above]{n\~ao} (r3.west);
  \draw[arrow] (ouro.south) -- node[lab,left]{sim} (persiste.north);

  \draw[arrow] (r3.south) -- node[lab,right]{sim} (descarta.north);
  \draw[arrow] (r3.west) to[out=180,in=0] node[lab,above,pos=0.3]{n\~ao} (persiste.east);
\end{tikzpicture}
 dentro de um ambiente figure.
\begin{tikzpicture}[
  font=\small,
  dec/.style={diamond, draw=black!60, fill=gray!8, align=center, inner sep=1pt, aspect=3, text width=2.2cm},
  act/.style={rectangle, draw=black!60, fill=green!6, rounded corners=2pt, align=center, inner sep=5pt, minimum width=4.0cm, text width=3.8cm},
  term/.style={rectangle, draw=black!60, fill=orange!12, rounded corners=2pt, align=center, inner sep=5pt, minimum width=4.2cm, text width=4.0cm},
  arrow/.style={-{Stealth[length=2.4mm]}, semithick},
  lab/.style={font=\scriptsize}
]
  \node[act] (inicio) at (0,3.0) {Leitura avaliada (rota\c{c}\~ao + picos $R^3$)};
  \node[dec] (primeira) at (0,0.6) {Primeira\\ leitura do\\ ponto?};
  \node[dec] (ouro) at (6.0,0.6) {$|\Delta RPM| > \Delta_{rot}$\\ e n\'umero de\\ picos $>0$?};
  \node[dec] (r3) at (12.0,0.6) {Todos os picos\\ dentro do\\ Paralelep\'ipedo?};
  \node[term] (persiste) at (6.0,-2.8) {PERSISTIR\\ \scriptsize leituras\_persistidas (nova refer\^encia)};
  \node[term] (descarta) at (12.0,-2.8) {DESCARTAR\\ \scriptsize registrar em leituras\_trash};

  \draw[arrow] (inicio) -- (primeira);

  \draw[arrow] (primeira.east) -- node[lab,above]{n\~ao} (ouro.west);
  \draw[arrow] (primeira.south) -- ++(0,-1.6) -| node[lab,above,pos=0.25]{sim} (persiste.west);

  \draw[arrow] (ouro.east) -- node[lab,above]{n\~ao} (r3.west);
  \draw[arrow] (ouro.south) -- node[lab,left]{sim} (persiste.north);

  \draw[arrow] (r3.south) -- node[lab,right]{sim} (descarta.north);
  \draw[arrow] (r3.west) to[out=180,in=0] node[lab,above,pos=0.3]{n\~ao} (persiste.east);
\end{tikzpicture}

\caption{Fluxo de decisão do motor de descarte: regra de ouro e Paralelepípedo de Descarte \(R^3\).}
\label{fig:fluxo-descarte}
\caption*{\footnotesize Fonte: elaboração própria, conforme `discard_engine.py`.}
\end{figure}

\begin{figure}[htbp]
\centering
% Diagrama: Máquina de estados por leitura simulada
% Uso: % Diagrama: Máquina de estados por leitura simulada
% Uso: \input{diagrams/estados} dentro de um ambiente figure.
\begin{tikzpicture}[
  font=\small,
  st/.style={rectangle, draw=black!60, fill=blue!6, rounded corners=3pt, align=center, inner sep=6pt, minimum width=3.0cm},
  fin/.style={rectangle, draw=black!60, fill=orange!12, rounded corners=3pt, align=center, inner sep=6pt, minimum width=3.0cm},
  arrow/.style={-{Stealth[length=2.4mm]}, semithick},
  lab/.style={font=\scriptsize}
]
  \node[st] (conf) at (0,0) {Configurado};
  \node[st] (sinal) at (4.0,0) {Sinal gerado};
  \node[st] (fft) at (8.0,0) {FFT/RMS/DC};
  \node[st] (aval) at (12.0,0) {Avaliando descarte};
  \node[fin] (pers) at (16.4,0) {Persistida};
  \node[fin] (desc) at (14.2,-2.4) {Descartada};

  \draw[arrow] (conf) -- node[lab,above]{gerar sinal} (sinal);
  \draw[arrow] (sinal) -- node[lab,above]{calcular FFT} (fft);
  \draw[arrow] (fft) -- node[lab,above]{avaliar} (aval);
  \draw[arrow] (aval.east) -- node[lab,above,pos=0.4]{regra de ouro / fora de $R^3$ / primeira leitura} (pers.west);
  \draw[arrow] (aval.south) to[out=-90,in=180] node[lab,left,pos=0.25]{todos dentro de $R^3$} (desc.west);
\end{tikzpicture}
 dentro de um ambiente figure.
\begin{tikzpicture}[
  font=\small,
  st/.style={rectangle, draw=black!60, fill=blue!6, rounded corners=3pt, align=center, inner sep=6pt, minimum width=3.0cm},
  fin/.style={rectangle, draw=black!60, fill=orange!12, rounded corners=3pt, align=center, inner sep=6pt, minimum width=3.0cm},
  arrow/.style={-{Stealth[length=2.4mm]}, semithick},
  lab/.style={font=\scriptsize}
]
  \node[st] (conf) at (0,0) {Configurado};
  \node[st] (sinal) at (4.0,0) {Sinal gerado};
  \node[st] (fft) at (8.0,0) {FFT/RMS/DC};
  \node[st] (aval) at (12.0,0) {Avaliando descarte};
  \node[fin] (pers) at (16.4,0) {Persistida};
  \node[fin] (desc) at (14.2,-2.4) {Descartada};

  \draw[arrow] (conf) -- node[lab,above]{gerar sinal} (sinal);
  \draw[arrow] (sinal) -- node[lab,above]{calcular FFT} (fft);
  \draw[arrow] (fft) -- node[lab,above]{avaliar} (aval);
  \draw[arrow] (aval.east) -- node[lab,above,pos=0.4]{regra de ouro / fora de $R^3$ / primeira leitura} (pers.west);
  \draw[arrow] (aval.south) to[out=-90,in=180] node[lab,left,pos=0.25]{todos dentro de $R^3$} (desc.west);
\end{tikzpicture}

\caption{Máquina de estados do processamento de uma leitura simulada.}
\label{fig:estados}
\caption*{\footnotesize Fonte: elaboração própria, conforme arquitetura do projeto.}
\end{figure}

\section{Requisitos e rastreabilidade}

A especificação formaliza 19 requisitos funcionais (FR-001 a FR-019) e 7 não
funcionais (NFR-001 a NFR-007), detalhados em `SPECIFICATION.md`. A
Tabela~\ref{tab:requisitos} agrupa os requisitos por módulo de implementação.

\begin{table}[htbp]
\centering
\caption{Agrupamento dos requisitos por módulo implementado.}
\label{tab:requisitos}
\begin{tabular}{@{}p{4.4cm}p{8.6cm}@{}}
\toprule
\textbf{Módulo} & \textbf{Requisitos associados} \\
\midrule
`signal_generator` & FR-001, FR-002, FR-003 \\
`fft_processor` & FR-004, FR-005, FR-006, FR-019, NFR-001 \\
`discard_engine` & FR-007, FR-008, FR-009, FR-010, FR-011 \\
`persistence` & FR-011, FR-012, FR-013, NFR-003, NFR-004 \\
`api` & FR-016, FR-017, FR-018, NFR-002 \\
`frontend` & FR-014, FR-015, FR-016, FR-017, FR-018, FR-019 \\
\bottomrule
\end{tabular}
\caption*{\footnotesize Fonte: SPECIFICATION.md do projeto.}
\end{table}

\section{Geração de sinais sintéticos}

O módulo `signal_generator` implementa a síntese no domínio do tempo. Para cada
defeito do catálogo, um conjunto de \emph{componentes espectrais} define a
ordem (múltiplo de \(f_r\)), o peso relativo e, quando aplicável, a geração de
bandas laterais em \(\pm 1X\). O sinal é composto pela soma de senoides
\begin{equation}
\label{eq:sintese}
  s(t) = \sum_i A_i \sin(2\pi f_i t + \varphi_i) + n(t),
\end{equation}
em que \(f_i = ordem_i \cdot f_r\), \(A_i = peso_i \cdot severidade\) e
\(n(t)\) é ruído branco gaussiano com desvio padrão igual ao parâmetro de ruído
de fundo. A Tabela~\ref{tab:perfis} resume os perfis implementados
(Tabela~\ref{tab:catalogo} contém o catálogo completo).

\begin{table}[htbp]
\centering
\caption{Perfis espectrais implementados (ordens e pesos relativos) para defeitos representativos.}
\label{tab:perfis}
\begin{tabular}{@{}p{4.2cm}p{8.8cm}@{}}
\toprule
\textbf{Tipo} & \textbf{Componentes: ordem (peso)} \\
\midrule
sem\_defeito & 1X (0,20); 2X (0,03); 3X (0,015) \\
desbalanceamento & 1X (1,00); 2X (0,07); 3X (0,04) \\
desalinhamento\_angular & 1X (0,50); 2X (0,80); 3X (0,25); 1,5X (0,15); 2,5X (0,12); 4X (0,12) \\
desalinhamento\_paralelo & 1X (0,40); 2X (1,00); 4X (0,40); 6X (0,18); 8X (0,08) \\
rocamento & 1X (1,00); 2X (0,50); 3X (0,30); 4X (0,20); 5X (0,12); + sub-harmônicos se severo \\
mancal\_frouxo & 1X--10X com pesos 0,70 a 0,08; + fracionários 0,5X/1,5X/2,5X \\
acoplamento\_defeituoso & 1X (0,55); 2X (0,75); 3X (0,30); 4X (0,15) \\
rolamento\_bpfi & BPFI (1,00); 2$\times$BPFI (0,40); 3$\times$BPFI (0,20) $\pm$ sidebands 1X \\
\bottomrule
\end{tabular}
\caption*{\footnotesize Fonte: `DEFECT_PROFILES` e `_componentes_do_defeito` em `backend/app/domain/signal_generator.py`.}
\end{table}

Duas decisões de engenharia são relevantes. A primeira é a \emph{derivação da
frequência a partir da rotação}: todas as frequências são recalculadas como
\(f = ordem \cdot RPM/60\), nunca armazenadas como valores absolutos fixos, o
que mantém a coerência cinemática em qualquer regime. A segunda é o
\emph{determinismo da fase por ponto de medição}: a fase de cada componente é
gerada a partir de uma semente derivada do identificador do ponto (referência
de keyphasor estável), o que torna as leituras sucessivas de um mesmo ponto
comparáveis em fase --- condição necessária para a verificação \(R^3\) ---
enquanto o ruído permanece estocástico. O trecho a seguir ilustra o núcleo da
síntese:

\begin{lstlisting}[language=Python, caption={Núcleo da síntese de sinal em signal_generator.py.}, label={lst:gen}]
for componente in _componentes_do_defeito(tipo_defeito, severidade, rng_fase):
    freq = componente.ordem * freq_rotacao_hz      # ordem x (RPM/60)
    amplitude = componente.peso * severidade
    fase = rng_fase.uniform(0, 2*pi)
    sinal += amplitude * np.sin(2*pi*freq*t + fase)
    if componente.sidebands_1x:
        for f_lat in (freq - fr, freq + fr):       # bandas laterais em +-1X
            if f_lat > 0:
                sinal += 0.2*amplitude * np.sin(2*pi*f_lat*t + fase_lat)
sinal += rng_ruido.normal(0, ruido_efetivo, n)
\end{lstlisting}

Para os defeitos com razão subsíncrona dinâmica (\emph{oil whirl} e \emph{whirl
por atrito}), a razão é amostrada em tempo de geração dentro de uma faixa
especificada (\(0{,}39\)--\(0{,}48\) e \(0{,}30\)--\(0{,}70\), respectivamente),
reproduzindo a variabilidade desses fenômenos. Para defeitos não lineares
severo-dependentes (roçamento com severidade acima de \(0{,}6\) e mancal frouxo),
sub-harmônicos e fracionários são adicionados, e o piso de ruído é elevado
(roçamento e folga truncam a forma de onda), conforme a literatura
\cite{randall2011,scheffer2004}.

\section{Processamento espectral (FFT)}

O módulo `fft_processor` calcula a FFT real (via `numpy.fft.rfft`), extrai os
picos e as métricas. A estimativa de \(F_{max}\) de cada defeito é feita pela
maior ordem do perfil (com piso de ordem 10) multiplicada por uma margem de
bandas laterais de \(2{,}0\):
\begin{equation}
\label{eq:fmax}
  F_{max} = \max\left(ordem_{maxima}, 10\right)\times 2{,}0 \times f_r .
\end{equation}
A validação de Nyquist rejeita configurações com \(f_s \le 2F_{max}\) antes de
qualquer processamento (NFR-001). A detecção de picos considera máximos locais
da magnitude acima de um \emph{limiar} absoluto, definido como o limiar relativo
configurável multiplicado pela maior amplitude do espectro:
\begin{equation}
\label{eq:limiar}
  \lambda = \lambda_{rel} \cdot \max_k |X[k]|,
\end{equation}
com os picos ordenados por amplitude e limitados a um máximo configurável. Para
cada pico, registram-se frequência, amplitude (escala de pico, \(2|X|/N\)) e
fase (em graus), compondo os vetores \(R^3\). As métricas RMS total, RMS de
picos, RMS de ruído e valor DC são calculadas conforme a
Seção~\ref{sec:r3} do Capítulo~\ref{cap:fundamentacao}, e o sinal de tempo é
decimado para a visualização.

\section{Motor de descarte}

O módulo `discard_engine` implementa a regra de ouro e o Paralelepípedo de
Descarte (Equações~\eqref{eq:regra-ouro} e~\eqref{eq:descarte} do
Capítulo~\ref{cap:fundamentacao}). Sua interface recebe a leitura atual, a
última leitura efetivamente persistida do ponto (ou `None` na primeira) e as
tolerâncias configuráveis (desvio de rotação, \(\delta_f\), \(\delta_A\) e
\(\delta_\phi\)). O resultado é uma decisão enumerada:

\begin{lstlisting}[language=Python, caption={Decisões possíveis do motor de descarte.}, label={lst:decisoes}]
PERSISTIR_PRIMEIRA_LEITURA   # nao ha referencia (FR-011)
PERSISTIR_REGRA_DE_OURO      # |dRPM| > desvio e ha picos (FR-007)
PERSISTIR_FORA_TOLERANCIA    # algum pico fora do R^3 (FR-010)
DESCARTAR                    # todos os picos dentro do R^3 (FR-009)
\end{lstlisting}

A comparação de cada pico atual é feita contra o pico de referência mais
próximo em frequência; se não houver pico de referência ou se qualquer desvio
exceder a tolerância, a leitura é persistida. Como o motor é uma função pura
(entrada \(\rightarrow\) saída), todas as combinações são testáveis
unitariamente sem banco de dados, e a concorrência sobre o ponteiro de última
leitura é tratada na camada de persistência, dentro da mesma transação que grava
a nova leitura.

\section{Persistência e modelo de dados}

A camada `persistence` utiliza SQLAlchemy 2.x assíncrono e segue a hierarquia
Planta~$\rightarrow$~Área~$\rightarrow$~Máquina~$\rightarrow$~Ponto. O schema
lógico (Tabela~\ref{tab:schema}) materializa a arquitetura não normalizada
adotada (NFR-003): os picos \(R^3\) de cada leitura são armazenados como um
documento JSON em coluna `JSONB`, priorizando a velocidade de ingestão em
relação à conveniência de consulta, conforme as diretrizes de origem.

\begin{table}[htbp]
\centering
\caption{Visão lógica do modelo de dados persistido.}
\label{tab:schema}
\begin{tabular}{@{}p{5.2cm}p{7.8cm}@{}}
\toprule
\textbf{Relação} & \textbf{Conteúdo principal} \\
\midrule
plantas / areas / maquinas & hierarquia com chaves estrangeiras \\
pontos & nó de medição; `ultima\_leitura\_persistida\_id` (ponteiro crítico) \\
leituras\_persistidas & picos\_r3 (JSONB), rms\_total/ruido/picos, valor\_dc, rotacao, níveis, timestamp \\
leituras\_trash & leitura descartada + `motivo\_descarte` \\
snapshots\_defeito & par sensor + tipo de defeito + vínculo com a leitura \\
\bottomrule
\end{tabular}
\caption*{\footnotesize Fonte: `.specs/codebase/ARCHITECTURE.md` e migração Alembic do projeto.}
\end{table}

A primeira migração Alembic cria as oito relações da
Tabela~\ref{tab:schema}. O estado compartilhado crítico é o ponteiro
`pontos.ultima\_leitura\_persistida\_id`: a decisão de descarte lê esse ponteiro
e grava a nova leitura na \emph{mesma transação}, com lock de linha, evitando
corridas entre simulações concorrentes do mesmo ponto. Leituras aprovadas são
escritas em `leituras_persistidas` e atualizam o ponteiro; leituras descartadas
vão para `leituras_trash`, sem alterar a referência.

\section{API REST}

A API é exposta pelo FastAPI com validação de schema Pydantic v2. Os endpoints
implementados são resumidos na Tabela~\ref{tab:endpoints}; os contratos
detalhados (requisição e resposta) constam do Apêndice~\ref{ap:contratos}.

\begin{table}[htbp]
\centering
\caption{Endpoints REST implementados.}
\label{tab:endpoints}
\begin{tabular}{@{}p{4.6cm}p{4.4cm}p{4cm}@{}}
\toprule
\textbf{Endpoint} & \textbf{Descrição} & \textbf{Requisitos} \\
\midrule
POST /simulacoes & gera sinal, FFT e decide descarte & FR-001 a FR-019, NFR-001 \\
POST /snapshots & registra par sensor+defeito & FR-016 \\
GET /health & verificação de saúde do serviço & --- \\
\bottomrule
\end{tabular}
\caption*{\footnotesize Fonte: `backend/app/api/` do projeto.}
\end{table}

A resposta de `POST /simulacoes` contém `sinal_tempo` (decimado), `picos_r3`,
`rms_total`, `rms_ruido`, `rms_picos`, `valor_dc`, `rotacao`, `limiar_picos` e
`limiar_amplitude`, a decisão (`descartada` e `motivo_descarte`),
`tempo_processamento_ms` e `taxa_descarte_acumulada`. Erros são mapeados como
`422` (entrada inválida ou violação de Nyquist), `404` (ponto inexistente) e
`500` (falha interna de persistência, sem corromper o ponteiro de referência).

\section{Front-end}

O front-end é uma SPA em React com gráficos Recharts, organizada em
componentes: `PainelSimulacao` (formulário), `PainelSinal` (sinal e RMS),
`GraficoFFT` (barras de espectro), `IndicadoresSimulacao` (tempo e taxa de
descarte), `ChecklistValidacao`, `SnapshotDefeito` e o componente `Help`
(tooltips de ajuda). O formulário reúne os parâmetros de máquina e de
aquisição, incluindo um controle de limiar relativo de picos; a validação no
cliente espelha o critério de Nyquist e bloqueia a submissão de parâmetros
inválidos (FR-018, NFR-001).

O gráfico de FFT exibe cada barra com a frequência em hertz e a ordem
normalizada \(N\) (frequência dividida pela rotação em hertz), uma linha
horizontal com o limiar absoluto e, abaixo do gráfico, o valor da rotação em
hertz (FR-015). O registro de \emph{snapshot} associa o sensor ao tipo de
defeito simulado e persiste via `POST /snapshots` (FR-016). O client HTTP
tipado trata erros de rede e erros de API de forma explícita, exibindo
mensagens compreensíveis ao usuário.

\section{Infraestrutura e execução}

O ambiente de execução local é definido por Docker Compose com três serviços:
`db` (PostgreSQL 16 com volume persistente), `backend` (imagem Python que executa
`alembic upgrade head` e inicia o Uvicorn na porta 8000) e `frontend` (build de
produção servido por nginx na porta 80, com proxy de `/api/` para o back-end).
Os scripts `start.sh` e `stop.sh` automatizam, respectivamente, o setup (checagem
do Docker, cópia do `.env`, build de imagem se ausente, subida com espera de
health e semeadura de um ponto de demonstração) e a derrubada do ambiente
(inclusive com remoção opcional de volumes).

\section{Decisões de engenharia (ADRs)}

As principais decisões de arquitetura, registradas no projeto como ADRs, foram:

\begin{itemize}
  \item \textbf{JSONB para picos \(R^3\)} (NFR-003): banco relacional com colunas
  JSON para os vetores de picos, priorizando ingestão; consultas por pico
  individual exigem operadores JSON do PostgreSQL (reavaliar com indexação se
  necessário).
  \item \textbf{Processamento síncrono no MVP}: chamada request/response direta
  via FastAPI; sem fila de mensagens. Reavaliar se o volume concorrente (NFR-005)
  crescer.
  \item \textbf{Sem autenticação nesta fase} (NFR-006): uso local/interno, CORS
  restrito ao front de desenvolvimento; a validação Pydantic é a única barreira
  no MVP.
  \item \textbf{Determinismo de fase por ponto} e \textbf{limiar relativo
  configurável}: decisões de domínio que habilitam, respectivamente, a
  comparabilidade de fase no \(R^3\) e o descarte de ruído de fundo (FR-019).
\end{itemize}

\section{Observabilidade}

Os logs de simulação são estruturados por requisição (parâmetros, decisão de
descarte e tempo de processamento). A resposta de cada simulação inclui o tempo
de processamento (FR-017) e a taxa de descarte acumulada (NFR-004), calculada
pela camada de persistência como a proporção de leituras descartadas em relação
às avaliadas por ponto.

% Capítulo 6 — Experimentos e resultados
\chapter{Experimentos e Resultados}
% Conteúdo a escrever.

% Capítulo 7 — Discussão
\chapter{Discussão}
% Conteúdo a escrever.

% Capítulo 8 — Conclusão
\chapter{Conclusão}
% Conteúdo a escrever.


% =====================================================================
%  PÓS-TEXTUAIS
% =====================================================================
\backmatter
\bibliographystyle{plain}
\bibliography{references}

\appendix
% Apêndice A — Contratos de API
\chapter{Contratos de API}
% Conteúdo a escrever.

% Apêndice B — Exemplos de execução
\chapter{Exemplos de Execução}
% Conteúdo a escrever.


\end{document}
